\section{Conclusion}
\label{sec:conclusion}

\begin{table}[t]
\centering\small
\begin{tabular}{p{0.12\columnwidth}p{0.48\columnwidth}p{0.26\columnwidth}}
\toprule
ID & Hypothesis & Status (\splitname) \\
\midrule
H1 & hop-level credit beats trace-level RL & not supported \\
H1b & not only extra supervision (vs B2b) & supported \\
H1c & dense reward matters (vs A1) & not supported \\
H2 & paths are valid and auditable & partially supported \\
H3 & smaller ID$\rightarrow$OOD gap & inconclusive \\
\bottomrule
\end{tabular}
\caption{Hypothesis status, following the criteria of Section~\ref{sec:setup} (\R{SeedsHOneHoppo} seed(s) per method in this build).}
\label{tab:claims}
\end{table}

We studied whether hop-level credit helps multi-hop table question answering by casting each hop as a discrete entity-selection action trained with PPO, in an environment where encoder, candidates and answer head are shared across methods and budgets are matched.
Hop-level credit beat trace-level RL given the same path signal only at the end (H1b supported) and produced the most valid paths among the main methods, but it did not beat answer-only trace-level RL (H1 not supported), the dense reward did not improve accuracy over a terminal-only reward (H1c not supported), and full-table SFT remains the most accurate method in this regime. Path validity stays low (H2 partially supported); \hop\ loses less accuracy on held-out table clusters in relative terms, but this is not established statistically (H3 inconclusive). Across seeds and variants, accuracy tracks how much evidence the policy collects, and most errors start at the first hop.
Our setting is small (frozen small backbones, 3{,}000 training questions, heuristic gold paths), so the conclusions apply to this regime.
Natural next steps are a trained rather than heuristic path signal, larger backbones, and additional benchmarks (Appendix~\ref{app:cut}).
